\documentclass[10pt]{article}
\usepackage[T1]{fontenc}
\usepackage[utf8]{inputenc}
\usepackage[french]{babel}
\usepackage{amsmath,amssymb,mathrsfs,mathpazo}
\usepackage{geometry}
\geometry{a4paper,margin=2cm}
\usepackage{graphicx,array,tabularx,tabularray,colortbl,multirow,multicol,booktabs,enumitem,pifont,ulem,fancybox,diagbox,cancel}
\usepackage{pstricks,pstricks-add,pst-plot,pst-tree,pst-node,pst-text,pst-func,pst-math,pst-bspline,pst-3dplot,pst-eucl,pst-solides3d}
\let\pscalculate\undefined
\usepackage{pst-bezier}
\usepackage{tikz,pgfplots,tkz-tab}
\usetikzlibrary{arrows,arrows.meta,patterns,calc,positioning,trees,decorations.markings,shapes.geometric,shapes.misc,shapes.arrows,intersections,decorations.pathmorphing,decorations.pathreplacing,backgrounds,shadows,babel}
\pgfplotsset{compat=1.18}
\usepgfplotslibrary{fillbetween}
\usepackage[np]{numprint}
\usepackage{listings,makecell,pas-tableur,twoopt}
\let\barre\undefined
\usepackage[pstricks]{tablvar}
\definecolor{darkgreen}{RGB}{0,100,0}
\providecommand{\red}{\color{red}}
\providecommand{\blue}{\color{blue}}
\providecommand{\green}{\color{green}}
\providecommand{\black}{\color{black}}
\providecommand{\white}{\color{white}}
\providecommand{\gray}{\color{gray}}
\providecommand{\cyan}{\color{cyan}}
\providecommand{\magenta}{\color{magenta}}
\providecommand{\yellow}{\color{yellow}}
\providecommand{\violet}{\color{violet}}
\providecommand{\orange}{\color{orange}}
\providecommand{\euro}{\texteuro}
\providecommand{\EUR}{\texteuro}
\providecommand{\ieme}{\textsuperscript{e}}
\providecommand{\degres}{\ensuremath{^\circ}}
\providecommand{\degree}{\ensuremath{^\circ}}
\providecommand{\vv}[1]{\overrightarrow{#1}}
\providecommand{\Coord}[1]{#1}
\providecommand{\decosix}{}
\providecommand{\starredbullet}{\ensuremath{\bullet}}
\pagestyle{empty}
\setlength{\parindent}{0pt}
\usepackage{adjustbox}\DeclareUnicodeCharacter{2236}{:}\DeclareUnicodeCharacter{2190}{\ensuremath{\leftarrow}}
\providecommand{\R}{\mathbb{R}}
\providecommand{\N}{\mathbb{N}}
\providecommand{\Z}{\mathbb{Z}}
\providecommand{\Q}{\mathbb{Q}}
\providecommand{\D}{\mathbb{D}}
\providecommand{\C}{\mathbb{C}}
\providecommand{\e}{{\mathrm{e}}}
\providecommand{\vect}[1]{\overrightarrow{#1}}
\providecommand{\vectt}[1]{\overrightarrow{\mathrm{#1}}}
\providecommand{\vvt}[1]{\overrightarrow{\mathrm{#1}}}
\providecommand{\np}[2][]{{#2~#1}}
\providecommand{\vv}[1]{{\overrightarrow{#1}}}
\renewcommand{\d}{{\mathrm{d}}}
\providecommand{\degres}{{^\circ}}
\providecommand{\degre}{{^\circ}}
\providecommand{\blue}{\color{blue}}
\providecommand{\red}{\color{red}}
\providecommand{\enskip}{\quad}
\providecommand{\llbracket}{\lbrack\!\lbrack}
\providecommand{\rrbracket}{\rbrack\!\rbrack}
\providecommand{\euro}{{\text{€}}}
\providecommand{\no}{{\text{n°}}}
\providecommand{\interval}[2]{{[#1;#2]}}
\providecommand{\intervFO}[2]{{[#1;#2[}}
\providecommand{\dotfill}{\ldots}

\providecommand{\ds}{\displaystyle}
\providecommand{\textding}[1]{#1}
\providecommand{\Oij}{\(\mathrm{O};\vec{\imath},\vec{\jmath}\)}
\providecommand{\Oijk}{\(\mathrm{O};\vec{\imath},\vec{\jmath},\vec{k}\)}
\providecommand{\Ouv}{\(\mathrm{O};\vec{u},\vec{v}\)}
\providecommand{\intervFF}[2]{\ensuremath{[#1;#2]}}
\renewcommand{\euro}{\ifmmode\text{\texteuro}\else\texteuro\fi}
\DeclareUnicodeCharacter{00B0}{\ifmmode^{\circ}\else\textdegree\fi}
\DeclareUnicodeCharacter{20AC}{\ifmmode\text{\texteuro}\else\texteuro\fi}
% Énoncés : Éducation nationale. Transcription et corrigés : APMEP, Denis Vergès et contributeurs. Découpage et classement : CoopMaths, Sébastien Lozano, Éric Elter et contributeurs. Adaptation de la présentation : Maths Explorer.
% CC BY-SA : https://coopmaths.fr/www/about/
% https://coopmaths.fr/alea/static/bac/2024/tex/bac_2024_05_sujet1_ameriquenord_3_cor.tex
\begin{document}
\noindent\textit{APMEP / CoopMaths. CC BY-SA. Présentation adaptée par Maths Explorer.}\par\medskip
\medskip

Le but de cet exercice est d'étudier la fonction $f$ définie sur l'intervalle $]0~;~+\infty[$ par : 

\[f(x) = x\ln \left(x^2\right) -\dfrac 1x.\]

\smallskip

\textbf{Partie A : lectures graphiques}

\medskip

On a tracé ci-dessous la courbe représentative $(\mathcal{C}_f)$ de la fonction $f$, ainsi que la droite $(T)$, tangente à la courbe $(\mathcal{C}_f)$ au point A de coordonnées $(1~;~-1)$.

Cette tangente passe également par le point B$(0~;~-4)$.

\begin{center}
\psset{xunit=1.5cm,yunit=0.6cm,arrowsize=2pt 3}
\begin{pspicture*}(-0.4,-7.2)(4,7.1)
\multido{\n=1+1}{4}{\psline[linewidth=0.2pt](\n,-7)(\n,7)}
\multido{\n=-7+1}{15}{\psline[linewidth=0.2pt](0,\n,)(4,\n)}
\psaxes[linewidth=1.25pt,labelFontSize=\scriptstyle,showorigin=false]{->}(0,0)(-0.4,-7)(4,7)
\psplot[plotpoints=2000,linewidth=1.25pt,linecolor=red]{0.1}{4}{x dup mul ln x mul 1 x div sub}
\psplotTangent{1}{3}{x dup mul ln x mul 1 x div sub}
\uput[l](3,6){\red $\mathcal{C}_f$}\uput[dr](3,5){$T$}
\uput[dr](1,-1){A}\uput[dr](0,-4){B}
\psline[linewidth=1.75pt,ArrowInside=->](1,-1)(2,-1)(2,2)
\uput[d](1.5,-1){1}\uput[r](2,0.5){3}
\end{pspicture*}
\end{center}

\begin{enumerate}
\item 
Le coefficient directeur de la tangente au point d'abscisse 1 est le nombre dérivé $f'(1)$ ; on lit sur le graphique $f'(1) =  \dfrac31 = 3$.

L'ordonnée à l'origine est égale à $- 4$, donc l'équation réduite de la tangente $(T)$ est 

\[M(x~;~y) \in (T) \iff y = 3x - 4.\] 
\item 
$\bullet~$ il semble que $f$ est concave sur ]0~;~1[ ;

$\bullet~$ il semble que $f$ est convexe sur $]1~;~+ \infty[$.


Le point A semble être un point d'inflexion de la courbe $(\mathcal{C}_f)$.
\end{enumerate}

\bigskip

\textbf{Partie B : étude analytique}

\medskip

\begin{enumerate}
\item ~

$\bullet~~$On a $\displaystyle\lim_{x \to + \infty} \dfrac 1x = 0, \: \displaystyle\lim_{x \to + \infty} x = + \infty$ et $\displaystyle\lim_{x \to + \infty} \ln \left(x^2\right) = + \infty$ ; donc par produit de limites $\displaystyle\lim_{x \to + \infty} f(x) = + \infty$.

$\bullet~~$Avec $f(x) = 2x \ln x - \dfrac1x$, on sait que $\displaystyle\lim_{x \to 0} x \ln x = 0$ et $\displaystyle\lim_{x \to 0}\dfrac 1x = + \infty$, donc 

$\displaystyle\lim_{x \to 0}f(x) = - \infty$ (l'axe des ordonnées est asymptote verticale de $(\mathcal{C}_f)$ au voisinage de zéro.
\item On admet que la fonction $f$ est deux fois dérivable sur l'intervalle $]0~;~+\infty[$.
	\begin{enumerate}
		\item 
Les fonctions $x \longmapsto x$, \: $x \longmapsto \ln \left(x^2 \right)$ et $x \longmapsto - \dfrac 1x$ sont dérivables sur $]0~;~+ \infty[$ et l'on a :

$f'(x) = \ln \left(x^2 \right) + x \times \dfrac{2x}{x^2} + \dfrac{1}{x^2} = \ln \left(x^2 \right)+  2 +  \dfrac{1}{x^2}$ ou $f'(x) = 2\ln x + 2 + \dfrac{1}{x^2}$ sur $]0~;~+ \infty[$.
		\item 


En dérivant $f'(x)$ on obtient :

$f''(x) = \dfrac2x - \dfrac{2x}{x^4} = \dfrac2x - \dfrac{2}{x^3} = \dfrac{2x^2 - 2}{x^3} = \dfrac{2\left(x^2 - 1\right)}{x^3} = \dfrac{2(x + 1)(x - 1)}{x^3}$.
	\end{enumerate}
\item 
	\begin{enumerate}
		\item 
		Comme $x^3 > 0$ sur $]0~;~+ \infty[$, le signe de $f''(x)$ est celui de $(x + 1)(x - 1)$ et comme $x + 1 > 1 > 0$, le signe de $f''(x)$ est celui de $x - 1$.

Conclusion :

$\bullet~~$Sur $]0~;~1[,\: x < 1, \: f''(x) < 0$ : la fonction est concave ;

$\bullet~~$Sur $]1~;~+\infty[,\: x > 1, \: f''(x) > 0$ : la fonction est convexe ;

$\bullet~~$pour $x = 1$, la dérivée seconde s'annule en changeant de signe : le point A est le point d'inflexion de $(\mathcal{C}_f)$.
		\item 
Du signe de $f''(x)$ on en déduit les variations de $f'$ qui est décroissante sur ]0~;~1[ puis croissante sur $[1~;~+ \infty[$, donc $f'(1) = 2 + 1 = 3$ (vu à la question \textbf{1.}).

Comme 3 est le minimum de $f'$, on en déduit que $f'(x) > 0$ sur $]0~;~+ \infty[$ et par conséquent la fonction $f$ est  strictement croissante sur $\R_+^*$.
		

	\end{enumerate}
\item
	\begin{enumerate}
		\item 
		
Tableau de variations de $f$ :
		
\begin{center}
\psset{unit=1.75cm,arrowsize=2pt 3}
\begin{pspicture}(7,3.8)

\psframe(7,3.8)\psline(0,2)(7,2)\psline(0,3)(7,3)\psline(0,3.5)(7,3.5)\psline(1,0)(1,3.8)
\psline[linecolor=red](1.1,0)(1.1,3.5)\psline[linecolor=red](1.15,0)(1.15,3.5)
\uput[u](0.5,3.5){$x$} \uput[u](1.125,3.5){$0$} \uput[u](3,3.5){$1$} \uput[u](5,3.5){$2$} \uput[u](6.7,3.5){$+\infty$} 
\rput(0.5,3.25){$f''(x)$}\uput[u](2,3.1){$-$}\uput[u](3,3.1){$0$}\uput[u](4,3.1){$+$}\uput[u](6,3.1){$+$}\uput[u](3,2){3}
\rput(0.5,2.5){$f'(x)$}\psline{->}(1.2,2.8)(2.75,2.1) \psline{->}(3.3,2.1)(6.8,2.8)
\psline{->}(1.5,0.5)(6.5,1.5)\uput[u](1.4,0){$-\infty$}\uput[d](6.6,2){$+ \infty$}
\rput(0.5,1){$f$}\rput(3,0.85){$-1$}\rput(5,1.2){$\approx 2,27$}
\end{pspicture}
\end{center}

Sur l'intervalle ]1~;~2[ la fonction $f$ est continue (car dérivable) et strictement croissante avec $f(1) < 0$ et $f(2) > 0$ : d'après le théorème des valeurs intermédiaires il existe un nombre unique $\alpha \in ]1~;~2[$ tel que $f(\alpha) = 0$.

		\item 
		

La calculatrice donne $f(1,3) \approx - 0,09$ et $f(1,4) \approx 0,23$ donc $1,3 < \alpha < 1,4$, puis $f(1,32) \approx -  0,02$ et $f(1,33) \approx 0,007$, d'où $1,32 < \alpha < 1,33$ et enfin 

$f(1,327) \approx - 0,003$ et $f(1,328) \approx 0,0004$, donc $\alpha \approx 1,33$ au centième près.

On sait que $\alpha$ est solution de $f(x) = 0 \iff x \ln \left(x^2 \right) - \dfrac1x = 0 \iff $

$ x\ln \left(x^2 \right) = \dfrac{1}{x}  \iff \ln \left(x^2\right) = \dfrac{1}{x^2}$ et enfin par croissance de la fonction exponentielle :

$\exp\left(\ln \left(x^2\right)\right) = \exp\left(\dfrac{1}{x^2}\right) \iff x^2 = \exp\left(\dfrac{1}{x^2}\right)$.

$\alpha$ vérifie cette équation donc $\alpha^2 = \text{exp} \left(\dfrac{1}{\alpha^2}\right)$.
	\end{enumerate}
\end{enumerate}

\bigskip

\end{document}
